\chapter{Nonlocality and the boundary of reality}\label{ch:nonlocal}

\section{Bell, CHSH, Tsirelson}

Bell showed that no local hidden-variable theory can reproduce every correlation of quantum
mechanics~\cite{Bell1964}. In the CHSH form~\cite{CHSH1969}, local theories obey $|S|\le2$,
while quantum mechanics reaches the Tsirelson bound $|S|=2\sqrt2$~\cite{Tsirelson1980}.
Violations were seen by Aspect and collaborators~\cite{Aspect1982} and, with detection and
locality loopholes closed simultaneously, by three groups in 2015~\cite{Hensen2015,Giustina2015,Shalm2015}.
Entanglement has been distributed over 1{,}200~km from the Micius satellite~\cite{Yin2017}.

\section{The IAM account}

\emph{Non-Locality and the Boundary of Reality}~\cite{MahaffeyNonlocal} places nonlocal
correlations in the photon sector, which in the dual-sector model carries no gravitational
decoherence and hence no duration (Chapter~\ref{ch:time}). A null geodesic has $d\tau=0$, so
two photons created together remain, in IAM's reading, at zero duration from each other
until one is absorbed. Correlations of any spatial separation then carry no signal and need
no propagation. For entangled photons the paper predicts $|S|=2\sqrt2$ with no decay over any
distance. For entangled \emph{matter}, the gravitational channel of Chapter~\ref{ch:gravdec} dephases each body in its pointer basis and
multiplies the pair's coherence by $c(t)=1-D(t)$, with $D(t)=E_q(t/\tau_{\mathrm{IAM}})/e$ the assumed ramp of that chapter and
$\tau_{\mathrm{IAM}}$ from Eq.~\eqref{eq:tauIAM}. For pointer-basis dephasing the largest CHSH value is
\begin{equation}
 S_{\max}(t)=2\sqrt{1+c(t)^2}
\end{equation}
(Chapter~\ref{ch:entanglement}, Eq.~\eqref{eq:ent:smax}). \derived{} It stays above 2 for any remaining coherence and reaches 2 only when
the coherence is gone; with the settings that are optimal for the undecohered pair, $S=\sqrt2\,(1+c)$, which falls below 2 at $c<0.414$,
i.e.\ $D>0.586$. \calc{} The ramp itself is assumed. \conjecture

\begin{checkbox}[title={Independent check: the photon-to-matter distance ratio}]
The paper supports sector separation by comparing entanglement distances: photons over
1{,}200~km~\cite{Yin2017} against matter over ``approximately 1.3~m (trapped ions, 2019
[Hensen et al., 2015])'', a ratio of a million. The citation is to the Delft loophole-free
test, which entangled two electron spins in NV centres separated by \textbf{1.3~km}, not 1.3~m.
The correct ratio is therefore about $10^3$. Even that is set by engineering, fibre and
free-space losses, rather than by any matter-sector physics. The comparison does not support
the sector split.

The distinguishing prediction that remains is the one above. Entangled massive systems lose
Bell violation on the timescale $\tau_{\mathrm{IAM}}$, with the sigmoidal profile of
Chapter~\ref{ch:gravdec}, while entangled photons never do. For every massive system
entangled to date (ions, atoms, NV spins, superconducting circuits, and mechanical
oscillators near $10^{-12}$~kg~\cite{Riedinger2018,Kotler2021}), $\tau_{\mathrm{IAM}}$ at
physical densities is far longer than the experiment. No test has yet been possible.
\end{checkbox}

\section{The exponent $5/2$}

The same paper argues that one exponent, $n=5/2$, connects quantum decoherence to horizon thermodynamics, obtained bottom-up from decoherence and
top-down from horizon consistency, and that it appears in the cosmic information-production rate as $D(a)^{5/2}$ and in the electron mass as
$\alpha^{5/2}$. The book's own check does not support the cosmological half. With the full $\Lambda$CDM background, the source term
$D^{5/2}\Omega_mf/T_H$ gives $E(a)\propto\exp(0.76-0.87/a)$, not the $\exp(1-1/a)$ the law uses; the two coefficients cross unity at
$n\approx3$--$4$, bracketing the analytic $n=7/2$ (Chapter~\ref{ch:theory}, Section~\ref{sec:numerics}). \calc{} The exponent $5/2$ therefore
appears in the electron mass of Chapter~\ref{ch:particle} and in the decoherence argument, and not in the cosmology, which takes $7/2$. Whether the
two exponents are related is open. \openprob
